
\subsection*{Theory}
The theoretical foundations of our compiler workflow are built on two results that we present in this section. First, we introduce some notation. We use the following matrix norms to talk about approximation errors:
\begin{enumerate}
	\item $\normf{A} = \sqrt{\tr(A\dg A})$, the Frobenius norm. 
	\item $\normo{A} = \sigma_{\rm max}(A)$, the operator or spectral norm. 
	\item $\normt{A} = \tr(\sqrt{A\dg A}) = \sum_{i=1}^m \sigma_i(A)$, the Schatten 1-norm or trace norm. Here, $A$ is assumed to be a $m\times m$ square matrix.
\end{enumerate}
In addition, we will use the diamond norm for quantum channels,
\begin{align}
	\normd{\mathcal{E}} = \sup_X \{ \normt{(\mathcal{E}\otimes \mathcal{I})(X)} ~;~ \normt{X}\leq 1 \},
\end{align}
where $\mathcal{I}$ is an identity channel of the same dimension as $\mathcal{E}$. This norm induces a \emph{diamond distance} between channels, $d_\diamond(\mathcal{E}, \mathcal{E}') = \frac{1}{2}||\mathcal{E} - \mathcal{E}'||_{\diamond}$.

Let $V$ be a target unitary on $n$ qubits, and $U_i$ be a set of approximate compilations of $V$ provided by a compiler. The compilation workflow can have probabilistic elements, and therefore we think of $U_i$ as independent, identically distributed (\emph{i.i.d.}) samples from some distribution of approximate compilations. The statistical properties we will demand of the compilation workflow are:
\begin{enumerate}
\item The compilations should all be close to the ideal, in the sense,
    \begin{align}
	\normf{U_i-V} \leq \epsilon, \quad \forall i		
		\label{eq:ass1}
	\end{align}
	\item The compilations have bounded variance, in the sense, $\E{U_i} \equiv \E{U}$ is independent of $i$, and
	\begin{align}
		\E{\normf{U_i - \E{U}}^2} \leq \epsilon', \quad \forall i
		\label{eq:ass2}
	\end{align}
\end{enumerate}
The expectations in these expressions and in the following are over the distribution of compilations generated by the compiler that satisfy Conditions 1 and 2, \ie
$ \E{f(U)} = \int {\rm d}\mu_V(U) f(U), $
where ${\rm d}\mu_V(U)$ is the measure over $SU(2^n)$ defined by the compiler output that satisfies Conditions 1 and 2 when the target is $V$ (the expectation symbol $\E{}$ should have a subscript $V$ but we omit this for ease).

First, we present a restatement of a well-known lemma.

\noindent \emph{\bf Lemma 1:} \textit{\textbf{(Mixing Lemma, Campbell \& Hastings)}} \emph{Let $V$ and $U_i$ be defined as above. Let there be probabilities $\{p_i\}_{i=1}^M, (\sum_{i=1}^M p_i=1)$, such that $\normf{\sum_{i=1}^M p_i U_i - V} \leq O(\epsilon^2)$. Then, $d_{\diamond}(\mathcal{U}, \mathcal{V}) = O(\epsilon^2)$, where $\mathcal{V}[\rho] = V \rho V^{\dagger}$ is the action of the target unitary, and $\mathcal{U}[\rho] = \sum_{i=1}^M p_i U_i \rho U_i\dg$ is the channel defined by averaging over the approximate compilations according to weights $p_i$.}

This result follows immediately from the \emph{Mixing lemma} established by Campbell \cite{PhysRevA.95.042306} and Hastings \cite{hastings2016turning}, and the observation that the Frobenius norm upper bounds the spectral norm, and therefore $\normf{\sum_{i=1}^M p_i U_i - V} \leq O(\epsilon^2) \implies \normo{\sum_{i=1}^M p_i U_i - V} \leq O(\epsilon^2)$.

It is important that the condition in Lemma 1, referred to as the \emph{reduced weighted ensemble error} (ReWEE) condition in the main text, is stated in terms of the Frobenius norm, which is easily computable with simple matrix operations. This allows us to efficiently verify it, and moreover, to specify an efficient method for determining the optimal weights for the ensemble channel $\mathcal{U}$. To see this, we explicitly write out the quantity to minimize
\begin{align}
    & \min_{\{p_i\}}~ \normf{\sum_{i=1}^{M} p_i U_i - V}^2 \nn \\
    &= \min_{\{p_i\}}~ \tr\left((\sum_{i=1}^{M} p_i U_i - V)\dg (\sum_i p_i U_i - V)\right) \nn \\
    &= \min_{\{p_i\}}~ \sum_{i,j=1}^{M} p_i p_j \tr(U_i\dg U_j) - \sum_{i=1}^{M} p_i 2\Re\tr(U_i\dg V) \nn \\
    &= \min_{\p}~ \frac{1}{2} \p^{\sf T} H \p + f^{\sf T} \p, \qquad \text{s.t.}\quad  p_i\geq 0, ~~ E^{\sf T} \p = 1 \label{eq:quadprog}
\end{align}
where $\p = [p_1, p_2, ..., p_{M}]^{\sf T}, H_{ij} = 2\Re \tr(U_i\dg U_j)$, $f_i = -2\Re\tr(U_i\dg V)$, and $E = [1, 1, ..., 1]^{\sf T}$. In the last line we have written the minimization as a standard quadratic program with inequality constraints. The matrix $H$ is a Gram matrix (since $\tr(A\dg B)$ is an inner product) and thus, positive semidefinite. This, together with the fact that the constraints define a convex set, implies that the quadratic program is convex and therefore very tractable. 

While the ReWEE condition is operationally important, we can gain more insight into the properties required by the compiler producing $U_i$ by deriving an alternative \emph{sufficient} condition on the expected properties of the compiler output distribution. To do so, first we observe that Lemma 1 holds in expectation as well. That is, if there exist probabilities $\{p_i\}_{i=1}^M, (\sum_{i=1}^M p_i=1)$, such that $\E{\normf{\sum_{i=1}^M p_i U_i - V}} \leq O(\epsilon^2)$, then $\E{d_{\diamond}(\mathcal{U}, \mathcal{V})} = O(\epsilon^2)$.

\noindent \emph{\bf Lemma 2:} \emph{Let $V$ and $U_i$ be defined as above. Given probabilities $\{p_i\}_{i=1}^M$, ($\sum_{i=1}^M p_i=1$), a sufficient condition for achieving $\E{\normf{\sum_{i=1}^M p_i U_i - V}} \leq O(\epsilon^2)$ is the \emph{reduced bias condition}: $\normf{\E{U}-V} \leq O(\epsilon^2)$.}

The proof of this Lemma, which follows from an application of the \emph{bias-variance-covariance} decomposition, a well-known decomposition of generalization error of ensembles of estimators in machine learning \cite{Ueda_1996, Brown_20005}, is provided in the SI, Section 4.

Lemma 2 provides the key property that guides the development of the compiler workflow described in the next section.  Quadratic reduction in diamond distance of an ensemble of compilations to the target circuit can be achieved if the bias of the compiler output can be quadratically reduced. Reduction of this bias is the overarching motivation for the steps in our workflow.

As the proof of Lemma 2 makes clear, if the reduced bias condition is satisfied, then one can achieve the quadratically reduced diamond distance error even with a uniform ensembles -- \ie the optimization of weights $p_i$ using the quadratic program above is not necessary, and one can simply use a uniform ensemble over compiler outputs $U_i$. However, we can get better performance by using a weighted ensemble. Even if the reduced bias condition is met, determining weights through the quadratic program allows for smaller ensemble sizes (typically much smaller) than the ensemble size bound specified by Lemma 2, $M \geq \epsilon'/\epsilon^4$ (see proof in SI, Section 4). 


Lemmas 1 and 2 motivate the formation of ensemble channels constructed from approximate compilations with bounded error $\epsilon$ to achieve an output error that scales as $O(\epsilon^2)$, \emph{if} the compiler output satisfies the ReWEE condition or the reduced bias condition. Note that these conditions can be efficiently checked because they are stated in terms of an easily computable norm (Frobenius), whereas the resulting guarantee is in terms of the diamond norm. Similarly, the quadratic program that makes the ensemble efficient (in terms of number of unitaries that compose the channel) is tractable due to its convexity and the simple matrix computations required to form $H$ and $f$ in Eq. \eqref{eq:quadprog}.

Our final result establishes convergence criteria for implementing an ensemble quantum channel defined by $\{U_i, p_i\}_{i=1}^M$. 

\noindent \emph{\bf Lemma 3:} \emph{Let $V$ and $U_i$ be defined as above, and let these unitaries act on a register of $n$ qubits; \ie ${\rm dim}~ V = {\rm dim}~ U_i = 2^n \equiv d$. Define $\mathcal{V}[\rho] = V \rho V^{\dagger}$ as a target channel, and $\mathcal{U}[\rho] = \sum_{i=1}^M p_i U_i \rho U_i\dg$ as an optimized ensemble channel, such that $d_\diamond(\mathcal{V}, \mathcal{U})\leq O(\epsilon^2)$. Further, let $\hat{\mathcal{U}}[\rho] = \frac{1}{T}\sum_{i=1}^T U_{\sigma(i)} \rho U_{\sigma(i)}\dg$, with $\sigma(i)\in \{1,...,M\}$, be the empirical channel defined by sampling and averaging over $U_i$ according to the distribution $\{p_i\}$. For $\epsilon, \delta \in (0,1)$, $\pr\left\{ d_\diamond(\mathcal{V}, \hat{\mathcal{U}}) < O(\epsilon^2) \right\} > 1-\delta$, given
\begin{align}
    T \geq \frac{2v + \frac{2}{3}R \left(\epsilon^2/(2d^2)\right) }{\left(\epsilon^2/(2d^2)\right)^2}\log \frac{2d^2}{\delta}
    \label{eq:T_bound}
\end{align}
Here, $v \equiv \normo{\sum_{i=1}^M p_i (J_i - J_{\mathcal{U}})^2}, R \equiv \max_i \normo{J_i - J_{\mathcal{U}}}$, and $J_i (and J_{\mathcal{U}}$) are the $d^2$-dimensional Choi matrices associated to the channel $\mathcal{U}_i[\rho] = U_i \rho U_i^{\dagger}$ (and $\mathcal{U}$).}

The proof of this result is provided in the SI, Section 4, and is an application of the Bernstein matrix concentration inequality to an upper bound on the diamond norm provided by the difference in Choi matrices \cite{watrous_2018}. This lower bound on the number of samples required for convergence, Eq. \eqref{eq:T_bound}, strictly scales as $O(1/\epsilon^4)$, but in practice we have found that this is a vast overestimate because $v = O(\epsilon^2)$, and hence $T \geq O(1/\epsilon^2)$ suffices. The much faster convergence is consistent with the data shown in Fig. \ref{fig:td_convergences}, where we show convergence in trace distance of channel outputs.

\subsection*{Compiler workflow}
In this section we describe our full end-to-end compiler workflow, and its implementation for the demonstrations presented in the main text. Much of the implementation relies on tools integrated into BQSKit \cite{bqskit}, an open-source numerical optimization based compiler.
We can comfortably scale to algorithms with hundreds of qubits, and target both NISQ and FT gate sets. Figure \ref{fig:schematic} shows a detailed depiction of the workflow, which was outlined in Figure \ref{fig:resut_workflow} of the main text.


\begin{figure}[ht!]
    \centering
    \includegraphics[width=\linewidth]{figures/FullWorkflow_brev_two.png}
    \caption{Detailed overview of our compilation workflow. The input circuit is partitioned into blocks of width 4. At this circuit width, we can leverage numerical optimization techniques\cite{bqskit, ntro} to find more efficient circuit representations for our inputted circuit. We diversify this initial set of solutions and pass this ensemble to a quadratic program. The quadratic program gives us the corresponding sampling distribution, with which we calculate the \emph{weighted ensemble error} of each block ensemble. We can then select those blocks which quadratically reduce the error, and recombine them back into a full circuit ensemble.}
    \label{fig:schematic}
\end{figure}

\subsubsection*{Circuit partitioning}
The first step in the workflow is to partition the input circuit into non-overlapping sub-circuits acting on $w$ qubits, referred to as \emph{blocks} of width $w$. In BQSKit, there are 2 partitioners which trade off quality (minimal number of partitions) for speed. The \emph{ScanPartitioner} works by iteratively selecting the best block of $w$ qubits from left to right in a circuit. After scoring all possible $w$-qubit blocks starting from the current frontier (which divides partitioned from un-partitioned sections), it picks the highest-scoring block and moves the frontier past it. This process repeats until the entire circuit is partitioned. On the other hand, the \emph{QuickPartitioner} avoids this exhaustive search of all possible blocks, and greedily groups gates that appear in topological order. For our implementation, we utilize the ScanPartitioner for smaller circuits, and rely on a combination of both partitioners for larger circuits. For these circuits, we first divide into larger blocks with the QuickPartitioner, and split those larger blocks into blocks of width $w$ with the ScanPartitioner.

The choice of block width, $w$, presents a tradeoff. Typically, the larger the block width, the more options there are for optimizing resources (\eg reducing expensive gates), and synthesizing many diverse circuit approximations of the block in the next step. However, the computational cost of approximate synthesis increases with $w$, and therefore this variable should be chosen based on computational resources and available time. In all the demonstrations in this paper, we choose $w=4$. For smaller block widths, we found that the solutions generated were rarely diverse enough to quadratically reduce the weighted ensemble error for the output ensemble. 

\subsubsection*{Ensemble Generation}
Our Ensemble Generation step of the workflow generates a set of circuits that approximate the target circuit for a block to Frobenius error $\epsilon$. The process is split into two main components: Numerical Synthesis and Diversification. Both steps are performed independently for each block and are therefore trivially parallelized. 

This step has an important distinction between the NISQ and FT workflows. When targeting NISQ architectures, Numerical Synthesis generates a single set of circuits in the CNOT + $U3$ gateset. Each member of this set of circuits satisfies $\normf{U_i-V}\leq \epsilon$. When targeting FT architectures, Numerical Synthesis generates two parallel outputs. The first output matches that of the NISQ workflow: a set of circuits in the CNOT + $Rz$ gateset that satisfies $\normf{U_i-V}\leq \epsilon$.  The second output is a \emph{single circuit} in the CNOT + $Rz$ gateset that satisfies the tighter threshold $\normf{U_i-V}\leq \epsilon ^2$. 

As we describe in further detail in the Diversification section, the reason for this is to maximize our resource reduction and probability of satisfying the ReWEE condition. The first output can achieve better resource reduction, as it gives a (quadratically) larger error budget to the Numerical Synthesis step. The second output still gives us resource reduction, but the ability to satisfy the ReWEE condition is much improved. By considering both options, we are able to choose the ensemble with the fewest resources that also quadratically reduces the error, and thereby improve the quality of our final output. Importantly, both workflows generate the same output at the end of the Diversification step: An ensemble of Clifford + T circuits that satisfies $\normf{U_i-V}\leq \epsilon$.

\emph{Numerical Synthesis}: The Numerical Synthesis step of the workflow constructs circuits with continuous angles that have \emph{fewer} resources than our target circuit. Our implementation utilizes the synthesis tools in BQSKit \cite{bqskit} for this step. Given a target unitary $V$ and native gate set $\mathcal{G}$, BQSKit constructs a family of circuits $U_i$ composed of gates in $\mathcal{G}$ that approximate $V$ by using a heuristic-guided search \cite{leap, qsearch} to generate different parametrized circuit ansatze. Then, it leverages numerical optimization to \emph{instantiate} the ansatze with the optimized parameters that minimize the Frobenius error to the target unitary, $\normf{U_i-V}$. The search step can be biased towards prioritizing circuits that minimize expensive resources such as CNOT gates or arbitrary angle rotations. In our implementation, we terminate the search and optimization as soon as the circuit ansatze contain as many expensive gates as the original target circuit, at which point we output all circuits found satisfying $\normf{U_i-V}\leq \epsilon$. While the universal gate set that each circuit is synthesized to is arbitrary, in the demonstrations presented in this paper, we synthesized to the CNOT+$U3$ gateset for NISQ implementations and the Clifford+$R_z$ gateset for FT implementations. 

For synthesis of blocks for FT implementations, on top of the heuristic-guided search, each circuit output by BQSKit undergoes another pass through the \emph{Numerical T gate Reduction} pass (NTRO) tool ~\cite{ntro}. NTRO optimizes the circuit by reducing the number of non-Clifford elements, which maintaining approximation error. NTRO reframes the instantiation problem as the constrained numerical optimization problem
\begin{equation}
    \min_\theta ~ \mathcal{C}_{\rm non-Clifford}(U_i(\theta)), \quad \text{ s.t. } \quad \|U_i(\theta) - V\|_F \leq \epsilon, \nn
\end{equation} 
where we write the synthesized circuit as $U_i(\theta)$ to explicitly indicate the dependence on the $R_z$ rotation angles within the circuit. The cost function $\mathcal{C}_{\rm non-Clifford}(U_i(\theta))$ is a heuristic that estimates the number of non-Clifford gates needed approximate unitary $V$ to precision $\epsilon$. NTRO minimizes this cost function by applying strategies that attempt to approximate each $R_z(\theta)$ in the circuit with a Clifford, or close to Clifford, rotation.

As mentioned above, in the FT case, there are two parallel outputs for Numerical Synthesis. The first output BQSKit and NTRO are executed multiple times to generate an ensemble of circuit satisfying $\|U_i(\theta) - V\|_F \leq \epsilon$. For the second output the same tools are applied to generate \emph{one} circuit, satisfying $\|U_i(\theta) - V\|_F \leq \epsilon^2$.

\emph{Diversification}: The next step in our workflow applies a number of heuristics to attempt to increase the diversity of the approximation ensembles for each block output by the Numerical Synthesis step. The aim of these heuristics is to generate an ensemble of approximation circuits that satisfy the ReWEE condition, \ie $\normf{\sum_i p_i U_i - V}\leq O(\epsilon^2)$. 

In the NISQ compilation case, each approximate solution output by Numerical Synthesis consists of CNOTs and $U3$ single qubit rotations. Since we do not want to increase the number of CNOTs our diversification algorithm only modifies the $U3$ gates. Our heuristic is motivated by the observation that for each approximate circuit $U_i$, with Frobenius error, $\normf{U_i - V} = e_i \leq \epsilon$, we have an excess error budget $\epsilon - e_i$. Our goal is to ``spend'' this remaining error budget to spread apart the initial set of solutions. A description of our heuristic solution is the following:
\begin{enumerate}
    \item Let the number of $U3$ rotations in the circuit be $N_u$. We can define the error budget per $U3$ gate as $\epsilon_{u3} = \frac{\epsilon - e_i}{N_u}$
    \item For each $U3$, generate a set of perturbations of distance $O(\epsilon_{u3})$ by generating a set of perturbing Hamiltonians, $\mathcal{P}$, that take the form $P_{abc} = aX + bY + cZ$ where $||[a,b,c]||_2 = \epsilon_{u3}$. Importantly, each time an operator with coefficients $[a,b,c]$ is added to the set, one with the opposing coefficient $[-a,-b,-c]$ must also be added to ensure local convexity. The total number of perturbations generated is user-specified.
    \item Generate a set of perturbed unitaries for each $U3$ by post-multiplication: $U3 \cdot e^{i P_{abc}}$, for all perturbating Hamiltonians $P_{abc} \in \mathcal{P}$. We then rewrite each perturbed single-qubit unitary as an explicit $U3$ rotation.
    \item Finally, we confirm that $\normf{U_i-V} \leq \epsilon$ for each perturbed version of $U_i$.
\end{enumerate}

In the FT compilation case, each approximate solution output by the Numerical Synthesis step consists of Clifford and $R_z$ gates. Each $R_z$ rotation needs to be approximated by a set of Clifford+$T$ single qubit rotations, and the non-deterministic procedure for doing this also allows us to introduce diversity. While we can perturb each $R_z$ angle by the error budget to expand the circuit, Campbell \cite{PhysRevA.95.042306} presents a method to generate a set of $O(\epsilon)$ approximations for an axial rotation with a resulting error of $O(\epsilon^2)$. The method depends upon an oracle to find an $O(\epsilon)$ approximation, for which we use Ross and Sellinger's Gridsynth\cite{ross_sellinger} algorithm. Gridsynth outputs a series of Clifford and $T$ gates that approximate an angle to a given precision. This allows us to perform an $Rz$ gate fault-tolerantly without additional ancilla qubits. It is shown in Kuchlinov \emph{et al.} \cite{Kliuchnikov_2023} that we can expand this methodology to ancilla-based techniques as well, which may save even more $T$ gates. This method not only diversifies our set of solutions, but removes additional $T$ gates in our final ensemble. At the end of this step, we again confirm that $\normf{U_i-V} \leq \epsilon$ for each perturbed version of $U_i$.

For the FT workflow, Diversification converts the output of synthesis to the target gate set, Clifford+$T$. Additionally, we are able to able to reduce the number of $T$ gates by compiling each $Rz$ to $O(\epsilon)$. Because of this, we are able to able to consider an additional workflow with a quadratically smaller budget passed to Numerical Synthesis ($\epsilon^2$) and a Diversification budget of $O(\epsilon)$. This ensemble has a much higher probability of satisfying the ReWEE condition, since we diversify each $Rz$ gate in a way that guarantees quadratic error reduction locally (with respect to the single qubit unitary). By considering a single circuit that has a distance $\epsilon^2$, the resultant ensemble is likely to form a channel with $O(\epsilon^2)$ distance from the target channel. This increased probability of meeting the ReWEE condition in the FT setting is reflected in the statistics of weighted ensemble error shown in SI, Section 3 (bottom plot of Figure 2). As mentioned, the trade off of this second workflow is that we limit the resource savings attained from the Numerical Synthesis step. By considering both outputs, we can choose the ensemble that saves the most resources and quadratically reduces the error.

In both the NISQ and FT cases, at the end of the diversification step, we end up with a larger ensemble for each block of partitioned circuit, with each member of the ensemble having $\epsilon$ error to the target unitary for that block.

\subsubsection*{Filtering out Blocks}
\label{sec:bias_filter}
The preceding Ensemble Generation steps are designed to construct ensembles of approximate circuits for each block of the circuit that satisfy the ReWEE condition. While this property is not guaranteed by these steps we can explicitly check it for each block. It is important to note that this verification is possible because (i) the optimization required to formulate the optimal weighted ensemble, and (ii) the ReWEE condition itself, are both efficiently computable. 

Therefore, in this step we perform the convex quadratic optimization over the ensemble for each block, and compute the weighted ensemble error, $\normf{\sum_{i=1}^{M^\k}p^\k_i U^\k_i - V^\k}$. If this quantity is $\leq C\epsilon^2$ we store this weighted ensemble for block $k$. In the demonstrations shown in the main text we chose $C=2$ for $\epsilon=10^{-1}$, and $C=20$ for $\epsilon \leq 10^{-2}$. If this error bound is not satisfied, we revert back to the original target circuit for the block, and the size of the ``ensemble'' for that block is one. In the FT workflow, we only revert back to the original circuit if both output ensembles do not satisfy the error bound. We refer to such blocks where the REWEE condition cannot be met while simultaneously reducing resources as ``rigid'' blocks. Intuitively, rigid blocks contain mostly fixed-angle gates, many of which are in the set of expensive gates. For example, if a block only contains a sequence of CNOT gates, it is impossible to reduce the number of CNOTs while also maintaining the ReWEE condition for $\epsilon$ below some value. In the SI we provide statistics that summarize the number of blocks that satisfy the ReWEE condition and the sufficient reduced bias condition for each of our benchmarks. 

$C$ is a hyper-parameter in this implementation of the workflow. Small values of $C$ provide tighter guarantees on the output error, while larger values allow for more resource savings. Note that even for large values of $C$, the overall output error can be bounded by $K\epsilon^2$ because even if some blocks have error close to $C\epsilon^2$, others can have much smaller error.

Note that there exist synthesis algorithms that guarantee that the ReWEE condition is met, \eg Campbell's convex hull algorithm \cite{PhysRevA.95.042306}. However, these provide no guarantees on resource usage and lead to extravagant use of expensive gates in order to satisfy the ReWEE condition for rigid blocks. Our technique maximizes generality and resource optimization while remaining scalable in order to be applied off-the-shelf to any input circuit. As a tradeoff, resource reduction is not possible for all blocks.

\subsubsection*{Generating the Full Circuit Ensemble}
The final step of our workflow is to combine all of the block ensembles into an ensemble over the entire circuit. 
Given optimized ensembles for $K$ partitions, this defines a joint distribution $\mathbb{P} \equiv \mathbb{P}_1 \times ... \times \mathbb{P}_K$, from which circuits are sampled when executing the full-scale circuit. 

